\originalpreface
本册围绕“一个部件的振动如何成为许多部件间传播的波, 波又如何形成光学图样”建立连续的物理课程. 读者应已学过牛顿定律、能量、电容电感、基本电磁场, 会求导、积分与二阶常系数微分方程.  电动力学和量子力学的完整理论各有独立册, 这里保留波动与光学需要的物理接口及明确使用条件.

\textbf{课程与教师.} 主线是 MIT 物理系本科课程 8.03SC Physics III: Vibrations and Waves, 2016年秋季, 教师 Yen-Jie Lee. OCW Scholar 版本提供24讲视频, 第1--23讲配有文字讲义; 第24讲为期末复习, 没有配套文字笔记. 本册依据已取得的23讲 Typed Notes 重编, 合并重复复习, 补齐推导, 校正有据的原件笔误. 2016是课程门户标注的学期, 不自动等于每份练习试卷的实际出题年份.

\textbf{先修与范围.} 物理先修采用8.01力学、8.02电磁学对应知识. 数学先修是微积分、向量和线性方程组; 谱定理与傅里叶完备性在应用时说明. 本轮没有选进阶光学课程或整本教材为重编范围. Howard Georgi, \emph{The Physics of Waves}（Benjamin Cummings, 1992）是官方指定教材; A. P. French 的同名主题教材及 Bekefi--Barrett 的电磁波教材为进一步阅读. 未取得的书题只列参考, 不补造题面. 

\textbf{大纲与课历.} 官方提供教学次序, 本册以讲次而非未公布的实际日历日期列课历. 第1--3讲为平衡、简谐、阻尼与驱动; 第4--6讲为正常模、拍及耦合驱动; 第7--9讲为对称、离散链、连续极限与驻波; 第10--11讲为行波、能量、反射和声波; 第12--16讲为电磁波、色散、傅里叶、波包、二维三维波和折射; 第17--19讲为偏振、波片、辐射和介质; 第20--22讲为薄膜、干涉、阵列、衍射与分辨率; 第23讲文字材料为量子波; 第24讲视频为复习. 精确讲次与本册章的对应见册末.

\textbf{考核.} 公开范围是10套作业（44道大题, 无官方作业答案）、Practice Exam 1与2、Practice Final Exam 1--3（共24道大题）及对应五份官方答案, 另有三张公式表. 本册把完整中文题面与答案纳入同一PDF, 保留原题和小问编号; 作业解答标为AI独立解答, 练习考试解答标为官方答案中文重编, 自行补推或校正另说明. 期中练习范围按门户为第1--8讲与第9--16讲, 期末范围为第1--23讲. 不把这些练习题册称作全部真实期中、期末试卷. Scholar大纲没有发布校园课程成绩权重、考试实际日期和未公开试卷, 本册不虚构. 

\textbf{署名与修改.} 原材料来源: Yen-Jie Lee, MIT OpenCourseWare, 8.03SC, Fall 2016, \url{https://ocw.mit.edu/courses/8-03sc-physics-iii-vibrations-and-waves-fall-2016/}. 本稿是非官方中文翻译与重编, 增加教学解释、自绘图和独立题解, 修改顺序与公式笔误, 不代表MIT审核或认可. 除明确标注的第三方例外外, 本册及相应改编依 \href{https://creativecommons.org/licenses/by-nc-sa/4.0/}{CC BY-NC-SA 4.0} 发布: 保留MIT与作者署名, 仅非商业使用, 再分发及改编采用相同许可并说明修改. 无关原书照片和视频未打包进本册; 已选原件和逐对象许可判断另存可追溯清单.
